\Needspace{6\baselineskip}
\begin{center}\hypertarget{native-260-article187}{}\pdfbookmark[2]{187.}{article-187}187.\end{center}
We return from these general investigations to the particular problem for
which we undertook them. Before we can determine an orbit as accurately
as possible from more observations than are strictly necessary, we must
already have an approximate determination that agrees fairly closely with
all the given observations. The quantities sought in the problem are the
corrections still to be applied to these approximate elements to obtain the
closest possible agreement. Since we may suppose that these corrections
will be so \SourceLeaf{261}small that their squares and products may be
neglected, the resulting changes in the calculated geocentric positions of
the celestial body can be found using the differential formulas given in
Section II of Book One. The positions calculated from the corrected elements
we seek will therefore be expressed as linear functions of the corrections
to the elements. Comparing them with the observed positions according to
the principles set out above will yield the most probable values. These
operations are so simple that they need no further illustration. It is clear
that any number of observations may be used, however far apart they are.
--- The same method may also be used to correct the \emph{parabolic} orbits of
comets when a longer series of observations is available and the closest
possible agreement is required.
